\begin{lemma}[A closed exit pair extends to three exit legs]
Let $G=(V,E)$ be a finite loopless multigraph whose elements of $E$ are
the exact edge identities, let $T$ be a tree-sequence, and let $\varphi$
be a proper partial $k$-edge-coloring.  Suppose that $V(T)$ is elementary
under $\varphi$.  Let $a,b\in[k]$ be distinct, suppose that
$a\in\overline\varphi(T)$ and the singleton $\{a\}$ is closed on $T$,
and suppose that $b\notin\overline\varphi(T)$.

Assume that $P_1,P_2$ are support-disjoint $T$-exit paths in the
$a,b$ incidence graph, with respective outside ends $a_1,a_2$ and tree
ends $b_1,b_2$, where $a_1\ne a_2$ and $b_1\ne b_2$.  Then there are a
third outside end $a_3$, a third tree end $b_3$, a $T$-exit path $P_3$
from $a_3$ to $b_3$, and an exact edge $f_3$ such that $P_3$ is
support-disjoint from both $P_1$ and $P_2$, its two ends are respectively
distinct from both preceding outside and tree ends, and $f_3$ lies on
$P_3$, is incident with $b_3$, and is a $b$-colored boundary-edge identity
of $V(T)$.
\end{lemma}
\begin{proof}
Write $X=V(T)$ using \noderef{TreeSequenceVertexSet}, and let $C_i$ be
the full \noderef{TwoColorComponent} through the outside vertex-node
$\operatorname{inl}(a_i)$ of $P_i$, for $i=1,2$.  We first prove, rather
than assume, that each $a_i$ has degree exactly one in its full component.
By \noderef{TreeExitPath}, $a_i\notin X$, at least one of $a,b$ is missing
at $a_i$, and $P_i$ has length at least two.  Thus $P_i$ supplies at least
one edge-node adjacent to $a_i$.  If $a$ is missing there, every neighbor
of $a_i$ in \noderef{TwoColorIncidenceGraph} is an exact $b$-colored edge
incident with $a_i$.  Two distinct such edges would be adjacent in the
sense of \noderef{AdjacentEdges} and contradict
\noderef{IsProperEdgeColoring}.  Hence there is exactly one neighbor.  The
same argument with the colors reversed applies if $b$ is missing.  This
also agrees with the global degree bound
\noderef{TwoColorIncidenceDegreeAtMostTwo}.

The path $P_i$ supplies an exact edge-node in $C_i$, while the degree-one
node $\operatorname{inl}(a_i)$ and $a\ne b$ supply the leaf required by
\noderef{TwoColorPathComponent}.  Thus $C_i$ satisfies
\noderef{TwoColorNontrivialPathComponent}.  Applying
\noderef{TwoColorComponentSpanningPath} gives a simple path $R_i$ whose
support is exactly $C_i$.  The degree bound shows that no internal node of
$R_i$ has an extra neighbor.  An extra adjacency between the two ends
would leave no degree-one node, and any other extra adjacency at an end
would give degree at least three.  Therefore the full component is exactly
this finite simple path.  Since $a_i$ has degree one, it is one of its
ends.  Moreover, the simple path $P_i$ is the unique segment of $R_i$ from
$a_i$ to $b_i$; its \noderef{TreeExitPath} condition says that $b_i$ is
the first vertex of $X$ on that oriented component path.

As in the endpoint analysis below, both ends of every such nontrivial
component are vertex-nodes.  An edge-node representing a colored exact
edge $f$ has, by \noderef{TwoColorIncidenceGraph}, precisely the
vertex-node neighbors incident with $f$ in the sense of
\noderef{Incident}.  The two entries of its endpoint pair are distinct by
\noderef{LooplessMultigraph}, so it has degree two.  A vertex-node end
misses $a$ or $b$: otherwise \noderef{PresentColors} provides distinct
incident exact edges of both colors and hence two neighbors.  This last
alternative is equivalently expressed by \noderef{MissingColors}.

We split according to whether $C_1=C_2$.  Suppose first that they are
different.  They are disjoint by
\noderef{TwoColorComponentsDisjoint}.  Besides the already known outside
ends $a_1,a_2$, let $c_1,c_2$ be the respective other ends of $R_1,R_2$.
These four vertex ends are distinct.  Because
$b\notin\overline\varphi(T)$, the union definition
\noderef{TreeMissingColors} says that no vertex of $X$ misses $b$.  Hence
every one of these ends that lies in $X$ must miss $a$.  Elementariness
from \noderef{ColoredVertexSetElementary} permits at most one such end in
$X$.  Since $a_1,a_2$ are already outside, at least one of $c_1,c_2$ is
outside.  Relabel the two full components for this paragraph so that
$c_1\notin X$.

Orient $R_1$ from $c_1$ and stop at its first vertex $b_3$ in $X$, which
exists because $P_1$ ends at $b_1\in X$.  Let $P_3$ be this first-contact
portion and put $a_3=c_1$.  It is simple, its only vertex in $X$ is its
last vertex, and its initial vertex is a full-component end missing $a$ or
$b$.  Its two distinct vertex endpoints and the bipartite definition
\noderef{TwoColorIncidenceGraph} give length at least two.  Thus every
clause of \noderef{TreeExitPath} holds for $P_3$.

The paths $P_2$ and $P_3$ lie in disjoint full components.  The paths
$P_1$ and $P_3$ are first-contact portions taken from opposite ends of the
same simple component path.  If they met, their ordered path intervals
could meet only at a common first contact $b_1=b_3$.  Let $g_1,g_3$ be
their respective final exact edge-nodes.  They are distinct by simplicity.
Their preceding vertex-nodes lie outside $X$, while their common final
vertex lies in $X$.  From \noderef{Incident},
\noderef{LooplessMultigraph}, and \noderef{TwoColorIncidenceGraph}, each
is therefore an exact member of \noderef{BoundaryEdges}, colored either
$a$ or $b$.  The $\{a\}$-closed hypothesis
\noderef{TreeColorSetClosed}, through the exact filter
\noderef{BoundaryColorEdges}, excludes color $a$.  Thus both $g_1,g_3$
are distinct $b$-colored edges incident with their common vertex, contrary
to \noderef{AdjacentEdges} and properness.  Consequently $P_3$ is
support-disjoint from $P_1$ as well as $P_2$.  Component disjointness and
support disjointness respectively give
$a_3\ne a_1,a_2$ and $b_3\ne b_1,b_2$.

Now suppose $C_1=C_2=C$.  The distinct vertices $a_1,a_2$ are degree-one
nodes of the same nontrivial simple component path, so they are precisely
its two ends.  Since $a\in\overline\varphi(T)$, the definition
\noderef{TreeMissingColors} supplies a vertex $x\in X$ at which $a$ is
missing.  It is unique: if another vertex of $X$ missed $a$, the two
missing-color sets would not be disjoint, contrary to
\noderef{ColoredVertexSetElementary}.

The assumption $b\notin\overline\varphi(T)$ and the complement relation
between \noderef{MissingColors} and \noderef{PresentColors} show that
$b$ is present at $x$.  Choose the resulting exact $b$-colored edge $g$
incident with $x$.  It gives a neighbor of $\operatorname{inl}(x)$ in the
two-color incidence graph.  Because $a$ is missing at $x$, every such
neighbor has color $b$, and properness permits at most one.  Thus $x$ has
degree exactly one in its component $C_x$.  The edge-node $g$ makes
$C_x$ nontrivial, so the same
\noderef{TwoColorPathComponent},
\noderef{TwoColorNontrivialPathComponent}, degree-bound, and spanning-path
argument used above shows that $C_x$ is a finite nontrivial simple path
having $x$ as one end.

The component $C_x$ differs from $C$.  Otherwise the one simple path $C$
would have the three distinct degree-one nodes $a_1,a_2,x$: the first two
are distinct by hypothesis, and $x$ is in $X$ while both $a_i$ are outside.
A nontrivial simple path has only its two ends as degree-one nodes, a
contradiction.  Hence \noderef{TwoColorComponentsDisjoint} makes $C_x$
disjoint from $C$.

Let $y$ be the other end of the spanning path for $C_x$.  It is a vertex,
is distinct from $x$, and misses $a$ or $b$ by the endpoint argument.  It
cannot lie in $X$: no vertex of $X$ misses $b$, while a second vertex of
$X$ missing $a$ would contradict the uniqueness of $x$.  Thus $y$ is
outside.  Orient this component path from $y$, stop at its first vertex
$b_3$ of $X$ (at worst $x$), set $a_3=y$, and call the resulting prefix
$P_3$.  The same first-contact verification gives all clauses of
\noderef{TreeExitPath}.  Since its full component $C_x$ is disjoint from
$C$, its support and both endpoints are disjoint from those of $P_1,P_2$.

In either main case, let $f_3$ be the final exact edge-node of $P_3$.
The vertex immediately before it is outside $X$, and $b_3$ is inside.
The endpoint and incidence argument using \noderef{Incident} and
\noderef{LooplessMultigraph} puts $f_3$ in \noderef{BoundaryEdges}; its
adjacency in \noderef{TwoColorIncidenceGraph} makes its color $a$ or $b$.
Closedness excludes $a$, so the definition
\noderef{BoundaryColorEdges} places the individual identity $f_3$ in the
$b$-colored boundary.  By construction it lies on $P_3$ and is incident
with $b_3$.  This completes every asserted field of the third exit leg.
\end{proof}
